\section{Background and Meta-Theoretical Foundations}
\label{sec:background}

Constructing a scientifically rigorous maturity model for cloud-native observability requires synthesizing principles from two distinct disciplines: (i) the mathematical foundations of system observability in distributed computing, and (ii) the methodological meta-theories governing Design Science Research (DSR) maturity artifacts in Information Systems and Software Engineering. This section establishes both pillars to ground the design of the CN-OMM.

\subsection{From Control-Theoretic Observability to Cloud-Native Telemetry}
\label{subsec:kalman_to_cloud}

In classical linear dynamical systems theory, Kalman \cite{kalman_1960} defined \textit{observability} as the mathematical dual of controllability. Given a continuous-time state-space system governed by state vector $\mathbf{x}(t) \in \mathbb{R}^n$, input vector $\mathbf{u}(t) \in \mathbb{R}^p$, and external output measurement vector $\mathbf{y}(t) \in \mathbb{R}^q$:
\begin{equation}
    \dot{\mathbf{x}}(t) = \mathbf{A}\mathbf{x}(t) + \mathbf{B}\mathbf{u}(t), \quad \mathbf{y}(t) = \mathbf{C}\mathbf{x}(t) + \mathbf{D}\mathbf{u}(t)
    \label{eq:kalman_state}
\end{equation}
the system is completely observable on $[t_0, t_1]$ if and only if the observability matrix $\mathcal{O}_{\text{Kalman}} = [\mathbf{C}^\top, (\mathbf{C}\mathbf{A})^\top, \dots, (\mathbf{C}\mathbf{A}^{n-1})^\top]^\top$ has full column rank, $\text{rank}(\mathcal{O}_{\text{Kalman}}) = n$. Full rank guarantees that any latent internal state $\mathbf{x}(t_0)$ can be uniquely reconstructed from finite external output trajectories $\mathbf{y}(t)$.

When transposed to cloud-native microservice architectures \cite{sridharan_2018, majors_2022}, the system is no longer a fixed-dimension linear plant, but a time-varying directed call graph $\mathcal{G}(t) = (\mathcal{V}(t), \mathcal{E}(t))$ where vertices $\mathcal{V}(t)$ represent ephemeral Kubernetes Pods and edges $\mathcal{E}(t)$ represent synchronous (HTTP/gRPC) or asynchronous message invocations. Furthermore, the external measurement operator $\mathbf{C}$ is not gratis: emitting, serializing, and exporting telemetry signals consumes physical compute cycles, memory allocations, and network I/O on the very nodes hosting $\mathcal{V}(t)$. Thus, cloud-native observability maturity cannot be defined merely by the volume of $\mathbf{y}(t)$; it must characterize the architectural capability to reconstruct latent fault states $\mathbf{x}(t)$ rapidly and accurately while bounding the physical perturbation (the \textit{Probe Effect}) imposed by the measurement operator itself.

\subsection{Design Science Research and Maturity Model Meta-Theory}
\label{subsec:dsr_foundations}

In Software Engineering and Information Systems, maturity models represent normative design artifacts that describe an evolutionary path from an initial, ad-hoc state to a target state of institutional and technical excellence \cite{debruin_2005, mettler_2011, adekunle_2022}. However, early maturity models frequently copied the 5-level nomenclature of the Capability Maturity Model (CMM) without domain-specific theoretical grounding or empirical verification \cite{lasrado_2015, proenca_2016}.

To enforce methodological rigor, De Bruin et al.\ \cite{debruin_2005} structured maturity model engineering into six sequential lifecycle phases (*Scope, Design, Populate, Test, Deploy, Maintain*), demonstrating that a complete model must evolve from purely \textit{descriptive} (assessing ``as-is'' capabilities) to \textit{prescriptive} (guiding ``to-be'' architectural improvements via stage-gates) and \textit{comparative} (benchmarking across systems). Complementarily, Becker, Knackstedt, and P\"{o}ppelbu\ss\ \cite{becker_2009} translated Hevner et al.'s seven DSR guidelines \cite{hevner_2004} into eight mandatory requirements (\textbf{R1--R8}): comparison with existing models (R1), iterative development (R2), rigorous evaluation (R3), multi-methodological design (R4), problem relevance (R5), formal problem definition (R6), targeted presentation (R7), and scientific documentation (R8). In their audit of 51 published maturity models, Becker et al.\ revealed that fewer than $18\%$ satisfied these scientific criteria.

Building upon these foundations, P\"{o}ppelbu\ss\ and R\"{o}glinger \cite{poppelbuss_2011, roglinger_2012} established formal design principles requiring maturity models to specify mutually exclusive dimensions, unambiguous level definitions, and verifiable maturation criteria. Meanwhile, Helgesson, H\"{o}st, and Weyns \cite{helgesson_2012} categorized maturity model evaluations into three rigor tiers:
\begin{itemize}[leftmargin=*, nosep]
    \item \textbf{Type 1 (Offline / Author / Theoretical Grounding):} Analytical verification of internal consistency, ontological completeness, and systematic comparison against prior models.
    \item \textbf{Type 2 (Expert / Simulated Assessment):} Qualitative evaluation via practitioner interviews, Delphi panels, or subjective case study questionnaires.
    \item \textbf{Type 3 (Real-System Empirical Evaluation):} Objective, quantitative validation performed by instantiating maturity levels within an operational system and measuring empirical process outcomes under controlled conditions.
\end{itemize}

\subsection{Process Theory and Configurational Equifinality in Maturity Progression}
\label{subsec:process_configurational}

A critical theoretical contribution identified by Lasrado, Vatrapu, and Andersen \cite{lasrado_2015} is that maturity models must explicitly declare their underlying \textit{Process Theory of Change} (Van de Ven and Poole \cite{vandeven_1995}) and their \textit{Configurational Logic} (Fiss \cite{fiss_2011}; van Steenbergen et al.\ \cite{vansteenbergen_2010}). Traditional CMM-inspired models implicitly assume a purely linear \textit{Lifecycle} motor, where every dimension advances in lockstep and overall maturity is computed as a compensatory arithmetic average.

In cloud-native systems, however, maturity evolution is governed by the interplay of two process motors \cite{vandeven_1995}:
\begin{enumerate}[leftmargin=*, nosep]
    \item \textbf{Teleological Motor:} Engineering teams purposefully instrument workloads toward an explicit operational goal (e.g., meeting Service Level Objectives and minimizing Mean Time to Diagnose, $MTTDiag$).
    \item \textbf{Dialectical Motor:} Every increase in telemetry granularity triggers an opposing physical thesis/antithesis conflict between \textit{Diagnostic Visibility} and the \textit{Instrumentation Probe Effect} ($\Delta_{\text{CPU}}, \Delta_{\text{RAM}}$, and storage cardinality explosion). Progressing to a higher maturity stage requires an architectural synthesis (e.g., eBPF zero-code probes, tail-based sampling, and index-free log/trace correlation) that resolves this dialectical tension.
\end{enumerate}

Furthermore, grounded in set-theoretic configurational theory \cite{fiss_2011, lasrado_2015}, observability dimensions are \textbf{non-compensatory}. Certain architectural capabilities act as \textit{necessary core conditions}: for example, deploying high-volume distributed tracing ($D_1$) without semantic cross-pillar correlation and exclusive self-time isolation ($D_2$) creates a configurational imbalance that degrades fault localization. Our formalization in Section~\ref{sec:cn_omm} directly encodes these stage-gate dependencies.
